\begin{figure}[ht]
\centering
\begin{tikzpicture}[x=1cm, y=1cm, font=\small]

\node[anchor=south, font=\bfseries] at (7, 6.1) {Air buoyancy: the sample and the reference weights displace different air volumes};

% Left pan: low-density sample (large volume)
\draw[thick] (2.5, 4.5) -- (2.5, 4.0);
\draw[thick, fill=gray!12] (1.4, 3.2) rectangle (3.6, 4.0);
\node[font=\scriptsize] at (2.5, 3.6) {sample};
\node[anchor=north, font=\scriptsize] at (2.5, 3.15) {volume $V_s$, density $\rho_s$ (low)};

% buoyant force up-arrow on sample
\draw[->, thick, engblue] (2.5, 3.2) -- (2.5, 2.4) node[anchor=north, font=\scriptsize, engblue] {$F_b = \rho_a V_s g$};

% Right pan: dense reference weights (small volume, same mass)
\draw[thick] (11.5, 4.5) -- (11.5, 4.0);
\draw[thick, fill=gray!35] (11.05, 3.5) rectangle (11.95, 4.0);
\node[anchor=west, font=\scriptsize] at (12.05, 3.75) {reference weights};
\node[anchor=west, font=\scriptsize] at (12.05, 3.35) {volume $V_r < V_s$, density $\rho_r$ (high)};

% buoyant force up-arrow on reference (smaller)
\draw[->, thick, engred] (11.5, 3.5) -- (11.5, 2.9) node[anchor=north, font=\scriptsize, engred] {$\rho_a V_r g$};

% Beam
\draw[ultra thick, engblack] (2.5, 4.5) -- (11.5, 4.5);
\draw[thick, fill=engblack] (7, 4.5) -- (6.7, 5.2) -- (7.3, 5.2) -- cycle;
\draw[thick] (7, 5.2) -- (7, 5.7);

% Air field hint
\foreach \x/\y in {0.8/2.2, 4.2/1.8, 5.3/2.6, 8.5/1.9, 9.6/2.5, 13.0/2.1} {
  \node[font=\scriptsize, enggray] at (\x, \y) {$\cdot$};
}
\node[anchor=west, font=\scriptsize, enggray] at (0.4, 1.3) {surrounding air, density $\rho_a \approx 1.2\,\text{g/L}$};

% Correction relation box
\node[anchor=north, font=\small] at (7, 1.0) {True mass $m = m_{\text{indicated}} \left[ 1 + \rho_a \left( \dfrac{1}{\rho_s} - \dfrac{1}{\rho_r} \right) \right]$};
\node[anchor=north, font=\scriptsize] at (7, 0.2) {The correction grows as the sample density $\rho_s$ falls below the reference density $\rho_r$; for a low-density sample it is not negligible.};

\end{tikzpicture}
\caption[Air-buoyancy correction in mass comparison]{A mass
comparison in air weighs the sample against reference weights whose
density differs from the sample's. Each body loses apparent weight by
the weight of the air it displaces, so a low-density sample of large
volume $V_s$ loses more than the compact reference of volume $V_r$,
and the balance reads the sample light. The correction restores the
in-vacuum mass from the in-air indication and the two densities; it
is negligible for dense metals near the reference density and grows
as the sample density falls, reaching about $0.1\,\%$ for a sample
near $1000\,\si{\kilogram\per\meter\cubed}$ weighed against
steel-class weights. Source: editors' schematic after the OIML R 111
buoyancy-correction convention.}
\label{fig:vol01:ch07:buoyancy-correction}
\end{figure}
